
\setcounter{footnote}{31}

Let \(R_{\varnothing}\) denote the presheaf that assigns the empty set to
every \(S\in\operatorname{Ob}\mathscr C\); it is an initial object of
\(\widehat{\mathscr C}\).\footnote{Editor's note: part \SGARefLink{sga3:iv:lemma:4:4:5:item:i:01}{(i)} of Lemma \SGARefLink{sga3:iv:lemma:4:4:5}{4.4.5} has
been corrected, and the proofs of all three parts have been given in detail.}

\medskip
\noindent\SGARefTarget{sga3:iv:lemma:4:4:5}\textbf{Lemma 4.4.5.}
\begin{enumerate}
  \item[(i)]\SGARefTarget{sga3:iv:lemma:4:4:5:item:i:01} \textit{Suppose that
    \(R_{\varnothing}\notin J(S)\) for every
    \(S\in\operatorname{Ob}\mathscr C\). If \((X_i)\) is a family of
    separated presheaves, then the coproduct presheaf
    \(\coprod_i X_i\) is separated.}\footnote{Editor's note: in general, the
    coproduct of two sheaves \(F,G\) is not a sheaf. Indeed, let
    \(S_1,S_2\in\operatorname{Ob}\mathscr C\). Suppose that the coproduct
    \(S=S_1\coprod S_2\) exists in \(\mathscr C\), and that the fiber product
    \(S_1\times_S S_2\) is an initial object \(\varnothing\) of
    \(\mathscr C\) (cf. I, \SGARefLink{sga3:i:subsection:1:8}{1.8}). Let \(R\) be the sieve of \(S\) with basis
    \(\{S_1,S_2\}\). Then \((F\coprod G)(R)\) is the disjoint union of
    \(F(S)\coprod G(S)\) and the sets
    \(F(S_i)\times G(S_j)\), for \(i\ne j\); thus \(F\coprod G\) is not a
    sheaf in general. On the other hand, let \(\mathscr C\) be the category
    with a single object \(S\) and \(\operatorname{id}_S\) as its only
    morphism, equipped with the topology defined by
    \(J(S)=\{R_{\varnothing},S\}\). Then the only separated presheaves are
    \(R_{\varnothing}\) and \(X=\mathbf h_S\) (which is a sheaf), and
    \(X\coprod X\) is not separated.}

  \item[(ii)]\SGARefTarget{sga3:iv:lemma:4:4:5:item:ii:01} \textit{Consider an equivalence relation in the category of
    presheaves}
    \[
      X\overset{u}{\underset{v}{\rightrightarrows}}Y,
    \]
    \textit{and let \(w:Y\to Z\) be the quotient. If \(X\) is a sheaf and
    \(Y\) is separated, then \(Z\) is separated.}

  \item[(iii)]\SGARefTarget{sga3:iv:lemma:4:4:5:item:iii:01} \textit{Consider an amalgamated sum in the category of
    presheaves, in which \(u\) and \(u'\) are monomorphisms:}
    \[
    \begin{array}{ccc}
      X & \xrightarrow{\ u\ } & Y\\
      {\scriptstyle u'}\big\downarrow && \big\downarrow\\
      Y' & \longrightarrow & V.
    \end{array}
    \]
    \textit{If \(Y\) and \(Y'\) are separated and \(X\) is a sheaf, then
    \(V\) is separated.}
\end{enumerate}

\noindent\textit{Proof.}
For \SGARefLink{sga3:iv:lemma:4:4:5:item:i:01}{(i)}, set \(X=\coprod_i X_i\). Let
\(S\in\operatorname{Ob}\mathscr C\), let
\(\tau:R\hookrightarrow S\) be a refinement of \(S\), and let
\(x_1,x_2\in X(S)\) satisfy \(x_1\circ\tau=x_2\circ\tau\).
There are indices \(i,j\) such that \(x_1\in X_i(S)\) and
\(x_2\in X_j(S)\). Since \(R\ne R_{\varnothing}\), there is a morphism
\(\phi:T\to R\) with \(T\in\operatorname{Ob}\mathscr C\). Then
\[
  x_1\circ\tau\circ\phi=x_2\circ\tau\circ\phi.
\]
Since \(X(T)\) is the disjoint union of the \(X_k(T)\), it follows that
\(i=j\). Because \(X_i\) is separated and is a subobject of \(X\), the map
\[
  X_i(S)\longrightarrow X_i(R)\longrightarrow X(R)
\]
is injective. Hence \(x_1=x_2\), which proves that \(X\) is separated.

We next prove \SGARefLink{sga3:iv:lemma:4:4:5:item:iii:01}{(iii)}. Let \(i:Y\to V\) and \(j:Y'\to V\) be the canonical
morphisms, and let \(K\) be the kernel of the pair
\[
  Y\times Y'
  \overset{p}{\underset{q}{\rightrightarrows}}V,
  \qquad
  p=i\circ\operatorname{pr}_1,\quad
  q=j\circ\operatorname{pr}_2.
\]
For \(S\in\operatorname{Ob}\mathscr C\), the monomorphisms \(u,u'\) allow us
to identify \(X(S)\) with its image in \(Y(S)\), respectively in \(Y'(S)\).
Write \(Z(S)\), respectively \(Z'(S)\), for the complement. Then \(V(S)\)
identifies with the disjoint union of \(Z(S)\), \(Z'(S)\), and \(X(S)\).
It follows at once that
\[
  i(S):Y(S)\longrightarrow V(S),
  \qquad
  j(S):Y'(S)\longrightarrow V(S)
\]
are injective and that
\[
  (u\times u')(S):X(S)\longrightarrow K(S)
\]
is bijective. Consequently \(i\) and \(j\) are monomorphisms, and
\(u\times u':X\to K\) is an isomorphism.

Set \(U=Y\coprod Y'\). For a refinement
\(\tau:R\hookrightarrow S\), there is a commutative diagram
\[
\begin{array}{ccc}
  U(S)&\xrightarrow{\ U(\tau)\ }&U(R)\\
  \big\downarrow&&\big\downarrow\\
  V(S)&\xrightarrow{\ V(\tau)\ }&V(R).
\end{array}
\]
Let \(v_1,v_2\in V(S)\) have the same image in \(V(R)\). By the definition
of \(V\), they lift to elements \(y_1,y_2\in U(S)\); let \(z_1,z_2\) denote
their images in \(U(R)\). Then \(z_1,z_2\) have the same image in \(V(R)\).

Since \(i:Y\to V\) and \(j:Y'\to V\) are monomorphisms, the maps
\(Y(R)\to V(R)\) and \(Y'(R)\to V(R)\) are injective. Thus, because \(Y\)
and \(Y'\) are separated, if \(y_1,y_2\) both belong to \(Y(S)\), or both
belong to \(Y'(S)\), then \(y_1=y_2\). Otherwise we may suppose that
\(y_1\in Y(S)\) and \(y_2\in Y'(S)\), so that \(z_1\in Y(R)\) and
\(z_2\in Y'(R)\). Since \(i\circ z_1=j\circ z_2\), the morphism
\[
  z_1\boxtimes z_2:R\longrightarrow Y\times Y'
\]
factors through \(K=X\). Moreover, since \(X\) is a sheaf, the restriction
map identifies \(X(S)\) with \(X(R)\). Hence there is an \(x\in X(S)\) such
that
\[
  u(x)\circ\tau=z_1=y_1\circ\tau,
  \qquad
  u'(x)\circ\tau=z_2=y_2\circ\tau.
\]
Since \(Y\) and \(Y'\) are separated, \(u(x)=y_1\) and \(u'(x)=y_2\), and
therefore \(v_1=v_2\). This proves that \(V\) is separated.

Finally, we prove \SGARefLink{sga3:iv:lemma:4:4:5:item:ii:01}{(ii)}. First observe that the morphism of presheaves
\[
  X\xrightarrow{\,u\boxtimes v\,}K,
\]
where \(K\) is the kernel of the pair
\[
  w\circ\operatorname{pr}_i:Y\times Y\rightrightarrows Z,
\]
is an isomorphism. Indeed, for every \(T\in\operatorname{Ob}\mathscr C\),
\(X(T)\) is an equivalence relation on \(Y(T)\), and hence the diagram
\[
  X(T)\xrightarrow{\,u\boxtimes v\,}Y(T)\times Y(T)
  \overset{w\circ\operatorname{pr}_1}
    {\underset{w\circ\operatorname{pr}_2}{\rightrightarrows}}
  Z(T)
\]
is exact in the category of sets.

Now let \(S\in\operatorname{Ob}\mathscr C\), and let
\(g_1,g_2:S\to Z\) be two morphisms that agree on a refinement
\(\tau:R\hookrightarrow S\). By the construction of \(Z\),
\[
  Z(S)=Y(S)/X(S),
\]
so there are morphisms \(f_1,f_2:S\to Y\) such that \(wf_i=g_i\). We have
\(wf_1\tau=wf_2\tau\); by the preceding observation, there is a morphism
\(\phi:R\to X\) such that
\[
  u\phi=f_1\tau,\qquad v\phi=f_2\tau.
\]
Since \(X\) is a sheaf, there is a morphism \(\psi:S\to X\) with
\(\phi=\psi\tau\). Thus, in \(Y(R)\),
\[
  u\psi\tau=f_1\tau,\qquad v\psi\tau=f_2\tau.
\]
Because \(Y(S)\to Y(R)\) is injective, \(Y\) being separated, it follows
that \(u\psi=f_1\) and \(v\psi=f_2\), hence \(g_1=g_2\). Thus \(Z\) is
separated.

